% GENERATED FILE. Edit canonical lessons, then run npm run generate:books.
% canonical-source-hash: fnv1a64:9088711f7e29b800
% canonical-lessons: GE-C355-projekt, GE-C355-auftrag, GE-C355-zusammenarbeiten, GE-C355-zusammenarbeit, GE-C355-organisieren

\chapter{Ein Projekt}
\label{ch:ge-ein-projekt}
\hlchaptermodality{355}

\begin{chapteropening}
\textbf{By the end of this chapter:} \emph{I can talk about a project and an order, working together and organizing the work.}

Complete the last lesson of chapter 355: I can talk about a project and an order, working together and organizing the work.
\end{chapteropening}

\section[das Projekt]{das Projekt — the project}

\label{lesson:GE-C355-projekt}



\begin{quote}
Before the new one: say the German for to present, then the German for the presentation.
\end{quote}

\subsection*{You'll want to know: das Projekt}
\textbf{das Projekt} — \textquotedblleft{}the project\textquotedblright{}.

\begin{quote}
Unser Projekt muss bis Ende Mai fertig sein.
\end{quote}

\begin{quote}
\emph{Our project has to be finished by the end of May.}
\end{quote}

Plural: \emph{die Projekte}. The stress is on the second syllable: \emph{pro-JEKT}.

\begin{cousinweb}
\emph{Projekt} is a learned borrowing from Latin \emph{prōiectum} (EtymDB), from \emph{prōicere} \textquotedblleft{}to throw forth\textquotedblright{} (Lewis \& Short), built on \emph{iacere} \textquotedblleft{}to throw\textquotedblright{}: a plan thrown forward.
\end{cousinweb}

\begin{grammarlens}[title={What you've built}]
One more word for telling, explaining and describing at B1.
\end{grammarlens}

\subsection*{Your turn}
\begin{itemize}
\raggedright
  \item \emph{Say it:} \emph{das Projekt}
  \item \emph{Say it:} \emph{das Projekt}, once more
  \item \emph{Say it:} a sentence of your own with \emph{das Projekt}
  \item \emph{Recall:} say the German for to present, then the German for the presentation, then say \emph{das Projekt} again
\end{itemize}

\begin{tcolorbox}[breakable,colback=teal!4,colframe=teal!35!black,title={Before you move on}]
What does das Projekt mean? (\textquotedblleft{}The project\textquotedblright{}.) Say it once more.
\end{tcolorbox}
\section[der Auftrag]{der Auftrag — the order, the task, the commission}

\label{lesson:GE-C355-auftrag}



\begin{quote}
Before the new one: say the German for the presentation, then the German for the project.
\end{quote}

\subsection*{You'll want to know: der Auftrag}
\textbf{der Auftrag} — \textquotedblleft{}the order, the task, the commission\textquotedblright{}.

\begin{quote}
Die Firma hat einen großen Auftrag aus dem Ausland bekommen.
\end{quote}

\begin{quote}
\emph{The company has received a large order from abroad.}
\end{quote}

Plural: \emph{die Aufträge}. \emph{Im Auftrag von} — on behalf of.

\begin{cousinweb}
\emph{Auftrag} is the noun of \emph{auftragen} \textquotedblleft{}to charge someone with a task\textquotedblright{}, built on \emph{tragen} \textquotedblleft{}to carry\textquotedblright{}, from Proto-Germanic \emph{draganą} (EtymDB): a task laid on someone to carry.
\end{cousinweb}

\begin{grammarlens}[title={What you've built}]
One more word for telling, explaining and describing at B1.
\end{grammarlens}

\subsection*{Your turn}
\begin{itemize}
\raggedright
  \item \emph{Say it:} \emph{der Auftrag}
  \item \emph{Say it:} \emph{der Auftrag}, once more
  \item \emph{Say it:} a sentence of your own with \emph{der Auftrag}
  \item \emph{Recall:} say the German for the presentation, then the German for the project, then say \emph{der Auftrag} again
\end{itemize}

\begin{tcolorbox}[breakable,colback=teal!4,colframe=teal!35!black,title={Before you move on}]
What does der Auftrag mean? (\textquotedblleft{}The order, the task, the commission\textquotedblright{}.) Say it once more.
\end{tcolorbox}
\section[zusammenarbeiten]{zusammenarbeiten — to work together}

\label{lesson:GE-C355-zusammenarbeiten}



\begin{quote}
Before the new one: say the German for the project, then the German for the order.
\end{quote}

\subsection*{You'll want to know: zusammenarbeiten}
\textbf{zusammenarbeiten} — \textquotedblleft{}to work together\textquotedblright{}.

\begin{quote}
Wir arbeiten seit drei Jahren eng zusammen.
\end{quote}

\begin{quote}
\emph{We have been working closely together for three years.}
\end{quote}

Separable: \emph{wir arbeiten zusammen}, \emph{wir haben zusammengearbeitet}.

\begin{cousinweb}
\emph{zusammenarbeiten} is built from \emph{zusammen} \textquotedblleft{}together\textquotedblright{} and \emph{arbeiten} \textquotedblleft{}to work\textquotedblright{}, from \emph{Arbeit}, Proto-Germanic \emph{arbaidiz} \textquotedblleft{}labour\textquotedblright{} (EtymDB).
\end{cousinweb}

\begin{grammarlens}[title={What you've built}]
One more word for telling, explaining and describing at B1.
\end{grammarlens}

\subsection*{Your turn}
\begin{itemize}
\raggedright
  \item \emph{Say it:} \emph{zusammenarbeiten}
  \item \emph{Say it:} \emph{zusammenarbeiten}, once more
  \item \emph{Say it:} a sentence of your own with \emph{zusammenarbeiten}
  \item \emph{Recall:} say the German for the project, then the German for the order, then say \emph{zusammenarbeiten} again
\end{itemize}

\begin{tcolorbox}[breakable,colback=teal!4,colframe=teal!35!black,title={Before you move on}]
What does zusammenarbeiten mean? (\textquotedblleft{}To work together\textquotedblright{}.) Say it once more.
\end{tcolorbox}
\section[die Zusammenarbeit]{die Zusammenarbeit — the cooperation, the teamwork}

\label{lesson:GE-C355-zusammenarbeit}



\begin{quote}
Before the new one: say the German for the order, then the German for to work together.
\end{quote}

\subsection*{You'll want to know: die Zusammenarbeit}
\textbf{die Zusammenarbeit} — \textquotedblleft{}the cooperation, the teamwork\textquotedblright{}.

\begin{quote}
Danke für die gute Zusammenarbeit!
\end{quote}

\begin{quote}
\emph{Thank you for the good cooperation!}
\end{quote}

No plural in everyday use. \emph{In Zusammenarbeit mit} — in cooperation with.

\begin{cousinweb}
\emph{Zusammenarbeit} joins \emph{zusammen} \textquotedblleft{}together\textquotedblright{} to \emph{Arbeit} \textquotedblleft{}work\textquotedblright{}, Old High German \emph{arbeit}, Proto-Germanic \emph{arbaidiz} (EtymDB).
\end{cousinweb}

\begin{grammarlens}[title={What you've built}]
One more word for telling, explaining and describing at B1.
\end{grammarlens}

\subsection*{Your turn}
\begin{itemize}
\raggedright
  \item \emph{Say it:} \emph{die Zusammenarbeit}
  \item \emph{Say it:} \emph{die Zusammenarbeit}, once more
  \item \emph{Say it:} a sentence of your own with \emph{die Zusammenarbeit}
  \item \emph{Recall:} say the German for the order, then the German for to work together, then say \emph{die Zusammenarbeit} again
\end{itemize}

\begin{tcolorbox}[breakable,colback=teal!4,colframe=teal!35!black,title={Before you move on}]
What does die Zusammenarbeit mean? (\textquotedblleft{}The cooperation, the teamwork\textquotedblright{}.) Say it once more.
\end{tcolorbox}
\section[organisieren]{organisieren — to organize}

\label{lesson:GE-C355-organisieren}



\begin{quote}
Before the new one: say the German for to work together, then the German for the cooperation.
\end{quote}

\subsection*{You'll want to know: organisieren}
\textbf{organisieren} — \textquotedblleft{}to organize\textquotedblright{}.

\begin{quote}
Meine Kollegin organisiert das Fest für die ganze Firma.
\end{quote}

\begin{quote}
\emph{My colleague is organizing the party for the whole firm.}
\end{quote}

Perfekt: \emph{hat organisiert}, with no \emph{ge-}. \emph{Sich organisieren} — to get organized.

\begin{cousinweb}
\emph{organisieren} was borrowed from French \emph{organiser} (EtymDB), from Medieval Latin \emph{organizāre}, built on Latin \emph{organum} \textquotedblleft{}instrument, tool\textquotedblright{}, from Greek \emph{órganon}.
\end{cousinweb}

\begin{grammarlens}[title={What you've built}]
One more word for telling, explaining and describing at B1.
\end{grammarlens}

\subsection*{Your turn}
\begin{itemize}
\raggedright
  \item \emph{Say it:} \emph{organisieren}
  \item \emph{Say it:} \emph{organisieren}, once more
  \item \emph{Say it:} a sentence of your own with \emph{organisieren}
  \item \emph{Recall:} say the German for to work together, then the German for the cooperation, then say \emph{organisieren} again
\end{itemize}

\begin{tcolorbox}[breakable,colback=teal!4,colframe=teal!35!black,title={Before you move on}]
What does organisieren mean? (\textquotedblleft{}To organize\textquotedblright{}.) Say it once more.
\end{tcolorbox}
